\section{Exponential passage times and the strict-convexity argument}
\label{model:section}

We now restrict to independent exponential edge weights of mean one.
The argument in this part uses no differentiability conclusion. Its
starting point is a hypothetical flat face of the limit shape, and its
contradiction will come from improving local paths in two directions
inside that face.

In this part $\tau_e$ has density $e^{-t}$ on $[0,\infty)$, and
$T$, $\mu$, and $\Bcal$ refer to this rate-one law.

We regard an edge also as an interval, assigning a partial traversal
the corresponding fraction of its cost and interpolating vertex
potentials affinely. This permits cuts inside edges while leaving all
vertex-to-vertex passage times unchanged.

The quantitative input is sublinear control of passage-time errors,
uniformly over polynomially many possible endpoints. We state the forms
used below for the exponential law.

\begin{lemma}[Uniform metric estimates]\label{model:metric}
There are constants $C,c>0$ such that, for $x,y\in\Z^2$ and
$R=|x-y|_1\ge2$,
\begin{align}
 0\le \E T(x,y)-\mu(y-x)
     &\le C\sqrt{R\log(2+R)},\label{model:mean}\\
 \PP\bigl(|T(x,y)-\E T(x,y)|\ge t\sqrt R\bigr)
     &\le C e^{-ct},\qquad t\ge0.\label{model:concentration}
\end{align}
Consequently, for any fixed $p,C_0<\infty$ and $\alpha\in(1/2,1)$,
with probability tending to one all pairs in $[-n^p,n^p]^2\cap\Z^2$
at distance at most $C_0n$ satisfy
\[
 |T(x,y)-\mu(y-x)|\le n^\alpha.
\]
The expected maximum lower deficit over those pairs is at most
$C_{p,C_0}\sqrt n\log(2+n)$. The same conclusions hold for
metric-graph endpoints after adding an $O(\log n)$ term in expectation.
\end{lemma}

\begin{proof}
Equation~\eqref{model:mean} is the nonrandom-fluctuation estimate of
Damron and Kubota \cite[Proposition~1.1]{DK2016}. The two-sided
concentration theorem of Damron, Hanson, and Sosoe
\cite[Theorem~1.1]{DHS2014} gives the stronger scale
$\sqrt{R/\log R}$; enlarging constants gives
\eqref{model:concentration}. Its hypotheses hold since
$\PP(\tau_e=0)=0$ and $\E e^{a\tau_e}=(1-a)^{-1}<\infty$ for $a<1$.
The finitely many small separations can be bounded by a coordinate
path and the exponential tail of its cost.

There are polynomially many pairs in the displayed box. Since
$\sqrt n\log n=o(n^\alpha)$, a union bound in
\eqref{model:concentration} proves the simultaneous assertion.
For the lower deficit, use $\mu(y-x)\le\E T(x,y)$ and integrate
the bound $\min\{1,C n^{C_p}e^{-ct/\sqrt n}\}$.
For interior-edge endpoints, attach the endpoint of each incident
edge, at a cost bounded by twice the maximum edge weight in a
slightly larger box; changes in $\mu$ are bounded by a fixed constant.
The expected maximum of polynomially many exponential variables is
$O(\log n)$. This proves the remaining statements.
\end{proof}

Lemma~\ref{lem:fb-cheap-paths} will supply a linear geodesic-length bound
with exponentially small failure probability by counting self-avoiding
paths. Together these estimates give the required microscopic controls.

\paragraph{A hypothetical flat and the quantity to improve.}
Suppose $\partial\Bcal$ contains a nontrivial segment. On its supporting
line take the maximal face, with endpoints $a,b$, and let $h$ be the
linear functional equal to one there. Then $h\le\mu$, and $\mu=h$ on
the positive cone generated by the face. The endpoints are extreme
points of $\Bcal$ and are linearly independent.

Choose rational vectors $v,w$ in this cone close to $a,b$. For a small
rational $s>0$ put $A=v+sw$ and $B=w+sv$. These slightly inward
directions will carry the two road families. Tracking and input
locality are measured by
$\|x\|_*=|[v\ w]^{-1}x|_\infty$. Constants may depend on the fixed
flat. The tracking and support constants of the final construction,
however, must be independent of $s$, of all large auxiliary cutoffs,
and of the input road rule; this is what will allow a fixed rescaling.

The variational quantity is the infimum $b(n)$ of a common bound for
the expected excess cost per period of the two local roads, with costs
kept in their original lattice units. We will prove both
$0<c_0\le b(n)$ and $b(n)\le C_s n^{3/4}$. The upper bound provides
inputs to an improvement; the lower bound makes indefinite improvement
impossible.

The middle of the proof establishes that improvement. Local gauges
make endpoint corrections computable from nearby inputs. Physical
routing bounds the savings that selected local tunnels can offer, and
a thin-diagram argument converts this into capped crossing estimates.
An independent spare road then replaces rare large positive costs.
The resulting pair satisfies $b(Dn)\le\rho b(n)$ for a fixed integer
$D$ and $\rho<1$, contradicting the positive lower bound.
